%%
%% This is file `sample-sigconf.tex',
%% generated with the docstrip utility.
%%
%% The original source files were:
%%
%% samples.dtx  (with options: `sigconf')
%% 
%% IMPORTANT NOTICE:
%% 
%% For the copyright see the source file.
%% 
%% Any modified versions of this file must be renamed
%% with new filenames distinct from sample-sigconf.tex.
%% 
%% For distribution of the original source see the terms
%% for copying and modification in the file samples.dtx.
%% 
%% This generated file may be distributed as long as the
%% original source files, as listed above, are part of the
%% same distribution. (The sources need not necessarily be
%% in the same archive or directory.)
%%
%%
%% Commands for TeXCount
%TC:macro \cite [option:text,text]
%TC:macro \citep [option:text,text]
%TC:macro \citet [option:text,text]
%TC:envir table 0 1
%TC:envir table* 0 1
%TC:envir tabular [ignore] word
%TC:envir displaymath 0 word
%TC:envir math 0 word
%TC:envir comment 0 0
%%
%%
%% The first command in your LaTeX source must be the \documentclass
%% command.
%%
%% For submission and review of your manuscript please change the
%% command to \documentclass[manuscript, screen, review]{acmart}.
%%
%% When submitting camera ready or to TAPS, please change the command
%% to \documentclass[sigconf]{acmart} or whichever template is required
%% for your publication.
%%
%%
\documentclass[sigconf]{acmart}

\usepackage{array}
\usepackage{breqn}
\usepackage{booktabs,makecell}
\usepackage{enumitem}
\usepackage{tikz,pgfplots}
\usetikzlibrary{positioning}
\usetikzlibrary{arrows}
\usetikzlibrary{patterns}
\usetikzlibrary{shapes.geometric}
\usetikzlibrary{decorations.pathreplacing}
\usetikzlibrary{calligraphy}
\usetikzlibrary{hobby}
\usetikzlibrary{calc}
\usepackage{thm-restate}

\usepackage{pifont}
\newcommand{\cmark}{\ding{51}}
\newcommand{\xmark}{\ding{55}}

\newtheorem{theorem}{Theorem}[section]
\newtheorem{corollary}{Corollary}[theorem]
\newtheorem{lemma}[theorem]{Lemma}
\newtheorem{claim}[theorem]{Claim}

\iffalse
\usepackage{draftwatermark}
\SetWatermarkScale{0.2}
\SetWatermarkAngle{270}
\SetWatermarkText{Accepted to SIGMOD 2025}
\SetWatermarkHorCenter{0.95\paperwidth}
\SetWatermarkVerCenter{0.5\paperheight}
\fi

\graphicspath{{./figures/}}

%% Align figure captions to the left
%\captionsetup{justification=raggedright,singlelinecheck=false}

%%
%% \BibTeX command to typeset BibTeX logo in the docs
\AtBeginDocument{%
  \providecommand\BibTeX{{%
    Bib\TeX}}}

%% Rights management information.  This information is sent to you
%% when you complete the rights form.  These commands have SAMPLE
%% values in them; it is your responsibility as an author to replace
%% the commands and values with those provided to you when you
%% complete the rights form.
\setcopyright{acmlicensed}
\copyrightyear{2025}
\acmYear{2025}
\acmDOI{XXXXXXX.XXXXXXX}

%% These commands are for a PROCEEDINGS abstract or paper.
\acmConference[VLDB '25]{VLDB International Conference on Very Large Databases
(2025)}{June 22--27, 2025}{London, United Kingdom}
%%
%%  Uncomment \acmBooktitle if the title of the proceedings is different
%%  from ``Proceedings of ...''!
%%
%%\acmBooktitle{Woodstock '18: ACM Symposium on Neural Gaze Detection,
%%  June 03--05, 2018, Woodstock, NY}
\acmISBN{978-1-4503-XXXX-X/18/06}


%%
%% Submission ID.
%% Use this when submitting an article to a sponsored event. You'll
%% receive a unique submission ID from the organizers
%% of the event, and this ID should be used as the parameter to this command.
%%\acmSubmissionID{123-A56-BU3}

%%
%% For managing citations, it is recommended to use bibliography
%% files in BibTeX format.
%%
%% You can then either use BibTeX with the ACM-Reference-Format style,
%% or BibLaTeX with the acmnumeric or acmauthoryear sytles, that include
%% support for advanced citation of software artefact from the
%% biblatex-software package, also separately available on CTAN.
%%
%% Look at the sample-*-biblatex.tex files for templates showcasing
%% the biblatex styles.
%%

%%
%% The majority of ACM publications use numbered citations and
%% references.  The command \citestyle{authoryear} switches to the
%% "author year" style.
%%
%% If you are preparing content for an event
%% sponsored by ACM SIGGRAPH, you must use the "author year" style of
%% citations and references.
%% Uncommenting
%% the next command will enable that style.
%%\citestyle{acmauthoryear}


%%
%% end of the preamble, start of the body of the document source.
\begin{document}

%%
%% The "title" command has an optional parameter,
%% allowing the author to define a "short title" to be used in page headers.
\title{Sketcbook Upper Bounds}

%%
%% The "author" command and its associated commands are used to define
%% the authors and their affiliations.
%% Of note is the shared affiliation of the first two authors, and the
%% "authornote" and "authornotemark" commands
%% used to denote shared contribution to the research.


%%
%% By default, the full list of authors will be used in the page
%% headers. Often, this list is too long, and will overlap
%% other information printed in the page headers. This command allows
%% the author to define a more concise list
%% of authors' names for this purpose.
%\renewcommand{\shortauthors}{Eslami, Dayan.}

%%
%% The abstract is a short summary of the work to be presented in the
%% article.
\begin{abstract}
  
\end{abstract}

%%
%% The code below is generated by the tool at http://dl.acm.org/ccs.cfm.
%% Please copy and paste the code instead of the example below.
%%
\begin{CCSXML}
<ccs2012>
   <concept>
       <concept_id>10003752.10003809.10010055.10010056</concept_id>
       <concept_desc>Theory of computation~Bloom filters and hashing</concept_desc>
       <concept_significance>500</concept_significance>
   </concept>
   <concept>
       <concept_id>10002951.10002952.10002971.10003450.10010830</concept_id>
       <concept_desc>Information systems~Unidimensional range search</concept_desc>
       <concept_significance>500</concept_significance>
   </concept>
 </ccs2012>
\end{CCSXML}

\ccsdesc[500]{Theory of computation~Bloom filters and hashing}
\ccsdesc[500]{Information systems~Unidimensional range search}

%%
%% Keywords. The author(s) should pick words that accurately describe
%% the work being presented. Separate the keywords with commas.
\keywords{Range Filter, Dynamic Data Structure, Data Growth, Scalability.}

\received{16 April 2024}
\received[revised]{4 June 2024}
\received[accepted]{8 August 2024}

%%
%% This command processes the author and affiliation and title
%% information and builds the first part of the formatted document.
\maketitle


\section{Theoretical Analysis and Results}\label{sec:theoretical_analysis_and_results}

\newcommand{\POECMS}{\mathsf{PoE_{CMS}}}
\renewcommand{\Pr}[1]{{\mathrm{Pr}}\left[ #1 \right]}
\newcommand{\parentheses}[1]{\left(#1\right)}
\newcommand{\expectation}[1]{\mathsf{E}\left[ #1 \right]}
\newcommand{\variance}[1]{\textsf{Var}\left[ #1 \right]}

We define $\POECMS(d,w,N)$ to denote the extra error we incur in CMS by not knowing the $\ell_1$ norm of the stream in advance. Namely, we compare the error of Sketchbook with that of a CMS with $d$ arrays of length $w$ after it has processed a stream of $\ell_1$ norm $N$.

\begin{lemma}
	In the CMS sketch, let $\mathcal{E}_{w^*}$ denote the error we get when we employ $w^*$ as a size function. Then, during epoch $i$:
	$$\expectation{\mathcal{E}_{w^*}} \leq  \frac{1}{w^*_i} \cdot 	\sum_{j=0}^i  2^{i-j} \cdot  (N_{j+1} - N_j) \;,$$
	where $N_j$ denotes the $\ell_1$ norm of the stream at the start of epoch $j$. 
\end{lemma}


\begin{lemma}
	In the CS sketch, let $\mathcal{E}_{w^*}$ denote the error we get when we employ $w^*$ as a size function. Then, during epoch $i$:
	$$\variance{\mathcal{E}_{w^*}} \leq  \frac{1}{w^*_i} \cdot 	\sum_{j=0}^i  2^{i-j} \cdot  (F_{j+1} - F_j) \;,$$
	where $F_j$ denotes the squared  $\ell_2$ norm of the stream at the start of epoch $j$. 
\end{lemma}



\textcolor{olive}{this needs more work, because it does involve arguing about the probability of deviating from the expected value..}
\section{Proof of Theorem for Insertion-only Streams}

We will do the analysis assuming a general size function $w^*$ and then discuss how the analysis proceeds for specific values of  $w^*$. Note that both CS and CMS share the same general structure: they contain $d$ counter arrays of size $w$. 

Fix some size function $w^*$. Note that we perform doubling each time the function $w^*$ doubles. In the CMS case, $w^*$ is a function of $N$, the $\ell_1$ norm of the stream. In the CS case,  $w^*$ is a function of the $\ell_2$ norm of the stream. In either case, we can track the moments \emph{in time} in which the value of the function doubles. Formally, instead of thinking of $w^*$ as a function of the the $\ell_1$ or $\ell_2$ norms of the stream, we think of it as a function of time such that, at any point in time $t$, $w^*(t)$ is the number of counters the sketch uses at time $t$.

Now let $T_0$ be the first time the sketch is instantiated (i.e., this sketch has $d$ counters and $w^*(T_0))$. %For simplicity's sake, we can assume  $w^*(T_0))$ is a power of two, and let $k \geq 0$ be such that $w^*(T_0))=2^k$.  
We then define $T_i$ for $i\geq 1$ to be the first moment in time when the length of each array doubles, i.e., $w^*(T_{i}) = 2\cdot w^*(T_{i-1})$. Note that, in between times $T_{i}$ and $T_{i+1}$, the length of each array stays the same (i.e., it is equal to $w^*(T_{i})$). We call this period of time epoch $i$ and denote by $\mathcal{S}_{i}$ the \emph{active} sketch during this epoch (i.e., the one that has arrays of length $w^*(T_{i})$). According to our design, updates to the stream during Epoch $i$ are executed directly on $\mathcal{S}_{i}$. For the purposes of executing queries fast, we also aggregate all previous sketches into the active sketch. 	

Formally, let $(x,\Delta)$ be an update that occured in some epoch $0\leq j < i$ before the current epoch. In the aggregation step, this update operation is replicated on all the counters in $\mathcal{S}_i$ that $x$ could have hashed into. Namely, if $x$ hashed to an index $\alpha$ in epoch $j$ for one of the arrays, then, during epoch $i$, all of the $2^{j-1}$ counters $\alpha \circ \{0,1\}^{j-i}$ ingest the update $(x,\Delta)$, where $\circ$ denotes concatenation and $ \{0,1\}^{j-i}$ denotes all binary strings of length $j-1$. From the perspective of fixed counter, what this means is that the counter must also ingest previous updates from elements that don't technically hash into it.

Formally, let $\alpha$ be a counter in epoch $i$, and assume for simplicity that we only have one array in the sketch (i.e., $d=1$ and $h_1(x) = h(x)$). This counter ingests the following operations:

\begin{itemize}
	\item updates $(x,\Delta)$ from epoch $i$ and for which $h(x) = \alpha$. This event occurs with probability $1/w^*(T_i)$ since there are $w^*(T_i)$ counters in the array. For ease of analysis, we define $\mathcal{J}_{i,\alpha}(x)$ to be the indicator random variable that is $1$ if this event occurs, and $0$ otherwise. Then $\Pr{\mathcal{J}_{i,\alpha}(x) =1}= 1/w^*(T_i)$.
	\item updates $(x,\Delta)$ from epoch $i-1$ that, during epoch $i-1$, hashed into a counter that matches the leftmost $\log w^*(T_i) - 1$ bits of $\alpha$ in the binary representation (meaning, their current $h(x)$ value differs from $\alpha$ in at most the last $1$ bit of their binary representation). This event occurs with probability at most $2/w^*(T_i)$. We emphasize that this is merely an upper bound: the additional factor of $2$ accounts for the mistakes that the aggregation step makes, i.e., processing updates from previous epoch that would not have hashed into $\alpha$ in the current epoch.
	\item  in general, updates $(x,\Delta)$ from epoch $j <i $ and for which $h(x)$ differs from $\alpha$ only in the last $i-j$ bits of its binary representation. This event occurs with probability at most $2^{i-j}/w^*(T_i)$. Similarly as before, we define $\mathcal{J}_{j,\alpha}(x)$ to be the indicator random variable that is $1$ if this event occurs, and $0$ otherwise. Then  $\Pr{\mathcal{J}_{j,\alpha}(x) =1}\leq  2^{i-j}/w^*(T_i)$. Similarly as before, note that this is merely an upper bound: for $x$ such that $h(x) = \alpha$, the $\Pr{\mathcal{J}_{j,\alpha}(x) =1 | h(x) = \alpha}=1$ no matter when the update to $x$ was processed in the stream.
\end{itemize}

While updates from previous epochs have a higher chance of hashing into $\alpha$ than updates in epoch $i$, they are nevertheless bounded in number, i.e., during epoch $i$, we have $T_{i+1}-T_{i}$ insertions. In the following, we bound the effect this has on each counter. We first cluster all the insertions that occur during a given epoch into one big insertion per element. Formally, let $f_i(x)$ denote the total count of updates  of the form $(x,\Delta)$ that occur in the stream during epoch $i$. Note that $\sum_x f_j(x) = T_{j+1}-T_j$ for all $j\geq 0$.

\medskip\noindent
\textbf{Bounds for CMS}. In the case of CMS, each counter gathers the sum of frequencies of elements that hash into it. Namely, the random variable $\alpha$ can be expressed as:

$$ \alpha =\sum_{j=0}^i   \sum_x \mathcal{J}_{j,\alpha}(x) \cdot f_j(x)   \;.$$

\noindent
By linearity of expectation, we get that:

\begin{align*}
	E[\alpha] &= \sum_{j=0}^i   \sum_x \frac{2^{i-j}}{w^*(T_i)} \cdot f_j(x) = \frac{1}{w^*(T_i)} \cdot  \sum_{j=0}^i   \sum_x 2^{i-j} \cdot f_j(x)\\
&=	\frac{1}{w^*(T_i)} \cdot  \sum_{j=0}^i   2^{i-j}  \sum_x f_j(x) = \frac{1}{w^*(T_i)} \cdot  \sum_{j=0}^i  2^{i-j} \cdot  (T_{j+1} - T_j) \;.
\end{align*}

In other words, the increase in error we get because of aggregating the sketches is dominated by the term $\sum_{j=0}^i  2^{i-j} \cdot  (T_{j+1} - T_j)$. This is where the specific size function $w^*$ we employ has an effect, since it translates into bounds on $T_i$ (and specifically their relative sizes). We now consider various special cases and discuss their respective behavior:

\medskip\noindent
\textbf{Setting $w^*(N) =  N/\epsilon$ in CMS.} In this case, we start with the original sketch $\mathcal{S}_0$ with a capacity for $T_0$ elements. The size function doubles whenever the number of elements doubles. In this case $T_{j+1} = 2 \cdot T_j$ for all $j\geq 0$ and, in general, $T_j = 2^j \cdot T_0$ for all $j\geq 0$. The additional error term then becomes:
\begin{align*}
	\sum_{j=0}^i  2^{i-j} \cdot  (T_{j+1} - T_j) & =	\sum_{j=0}^i  2^{i-j} \cdot  (2^{j+1}\cdot T_0 - 2^{j}\cdot T_0)\\
	&=  \sum_{j=0}^i  2^{i-j} \cdot  2^j \cdot T_0 = i \cdot 2^i \cdot T_0  \;.
\end{align*}	
Since the $\ell_1$ norm of the stream doubles in each epoch, we have that $2^i \cdot T_0 = T_i$ and so, in expectation,
$$ E[\alpha]  =  \frac{1}{w^*(T_i)} \cdot i \cdot T_i = i/\epsilon\;.$$
This is also an upper bound on the expected error we get for an element $x$. \textcolor{olive}{blah blah}

\medskip\noindent
\textbf{Setting $w^*(N) =  N^{1-c}/\epsilon$ in CMS, for some $0<c<1$.} In this case, the size function doubles more rarely, whenever $T_{i+1} = 2^{1/(1-c)} \cdot T_i$. Since $0<c<1$, note that this implies that $T_{i+1}$ is bigger than $2T_i$. This difference is why we avoid the extra dependency on $i$ in the error term. Formally, we have that $T_i = 2^{i/(1-c)} \cdot T_0$ and thus:

\begin{align*}
	\sum_{j=0}^i  2^{i-j} \cdot  (T_{j+1} - T_j) & =	\sum_{j=0}^i  2^{i-j} \cdot  (2^{(j+1)/(1-c)}\cdot T_0 - 2^{j/(1-c)}\cdot T_0)\\
	&=  \sum_{j=0}^i  2^{i-j} \cdot  2^{j/(1-c)} \cdot T_0 \cdot (2^{1/(1-c)} -1) \\
	&= 2^i \cdot T_0 \cdot (2^{1/(1-c)} -1) \cdot   \sum_{j=0}^i   2^{j/(1-c) - j}   \;.
\end{align*}	

For simplicity, we define $2^{1/(1-c)-1} = a$ and observe that:
$$ \sum_{j=0}^i   2^{j/(1-c) - j}   =  \sum_{j=0}^i   a^j =  (a^{i+1}-1)/(a-1) \;.   $$

We then get that:

\begin{align*}
	\sum_{j=0}^i  2^{i-j} \cdot  (T_{j+1} - T_j) & = 2^i \cdot T_0 \cdot (2^{1/(1-c)} -1) \cdot   \sum_{j=0}^i   2^{j/(1-c) - j}  \\
&= 2^i \cdot T_0 \cdot (2a -1) \cdot (a^{i+1}-1)/(a-1)  \\
	&\leq  2^{i+1} \cdot T_0 \cdot (a^{i+1}-1) \leq  2^{i+1} \cdot T_0 \cdot a^{i+1} \;.\\
\end{align*}	

At this point, it is worth mentioning that the expected error we are comparing against is $T_i/w^*(T_i)$, i.e., the error that we get by aggregating is close to the error we would get if we didn't expand and just allocated $S_i$ as a sketch from the beginning. To this end, we first notice that $2a = 2^{1/(1-c)}$ and so
$$ 2^{i+1} \cdot T_0 \cdot a^{i+1} = 2^{(i+1)/(1-c)} \cdot T_0  = 2^{1/(1-c)} \cdot T_i\;.$$

\noindent Therefore, the expected error becomes:
\begin{align*}
	E[\alpha]  &\leq   \frac{1}{w^*(T_i)} \cdot 	\sum_{j=0}^i  2^{i-j} \cdot  (T_{j+1} - T_j) \leq  2^{1/(1-c)} \cdot \frac{T_i}{w^*(T_i)} \;,
\end{align*}
as desired. We note that the above holds whenever $c<1$. In particular, $c$ does not need to be a constant. 
%
%We can expand the last term and get that:
%\begin{align*}
%	 \sum_{j=0}^i  2^{i-j} \cdot  (T_{j+1} - T_j) & =T_{i+1}-2^i \cdot T_0 + \sum_{j=1}^{i+1} 2^{i-j} \cdot T_i
%\end{align*}	

\medskip\noindent
\textbf{Streams with deletions.} \textcolor{olive}{to be filled}

\medskip\noindent
\textbf{Bounds for CS.} In the case of CS, we are interested in describing the behavior of the random variable $\hat{f}(z) = \alpha \cdot \sigma(z)$ for any query $z$, where $\sigma(z)$ is the associated sign function (for simplicity, we again assume that $d=1$). Using the notation we introduced before, the random variable $\alpha$ can be expressed as:

$$ \alpha =\sum_{j=0}^i   \sum_x \mathcal{J}_{j,\alpha}(x) \cdot f_j(x) \cdot \sigma(x)  = \sum_x \mathcal{C}(x)  \;,$$

where  $\mathcal{C}(x) =  \sum_{j=0}^i   \mathcal{J}_{j,\alpha}(x) \cdot f_j(x) \cdot \sigma(x)$ denotes the contribution throughout the epoch of updates to $x$.

We begin by computing the expected value of the estimator $\hat{f}(z)$ as in the classical CS analysis that shows that in expectation, $\alpha \cdot \sigma(z)$ is an unbiased estimator, i.e. $E[\alpha \cdot \sigma(z)] = f(z)$. By linearity of expectation

$$	\textsf{E}[\hat{f}(z)] =  \sum_{x \neq z}\expectation{\mathcal{C}(x) \cdot \sigma(z) } + \expectation{\mathcal{C}(z)\sigma(z)} = \sum_{j=0}^i f_j(z) = f(z)\;,$$

where the last step uses the fact that $\expectation{\mathcal{C}(x) \cdot \sigma(z) } =0$ by construction and that $\Pr{ \mathcal{J}_{j,\alpha}(z)}=1$ since $h(z)=\alpha$. We will now bound the variance of this estimator, and in particular, we will be interested in how much bigger it becomes when compared to $\sqrt{F_i/w^*(T_i)}$, where $F_i$ is the squared $\ell_2$ norm of the stream at time $T_i$.  Using a similar argument when expanding the variance of the estimator, we get that:

\begin{align*}
	\variance{\hat{f}(z)} &= \expectation{\parentheses{\sum_x \mathcal{C}(x) \sigma(z)}^2}  - (f(z))^2 \\
	& = \sum_x \expectation{\parentheses{\mathcal{C}(x)}^2} + \sum_x \sum_{y \neq x}\expectation{ \mathcal{C}(x) \mathcal{C}(y)} -  (f(z))^2 \\
	&=\sum_x \expectation{\parentheses{\mathcal{C}(x)}^2} - (f(z))^2 \;,
\end{align*}
where we  use the fact that $\expectation{\sigma(x) \cdot \sigma(y) } =0$  for any $x\neq y$ to establish that $\expectation{ \mathcal{C}(x) \mathcal{C}(y)}=0$. We now bound $\expectation{\parentheses{\mathcal{C}(x)}^2}$. Recall that we defined $\mathcal{C}(x) =  \sum_{j=0}^i   \mathcal{J}_{j,\alpha}(x) \cdot f_j(x) \cdot \sigma(x)$. Then:

\begin{align*}
	\expectation{\parentheses{\mathcal{C}(x)}^2} &= \expectation{\parentheses{ \sum_{j=0}^i   \mathcal{J}_{j,\alpha}(x) \cdot f_j(x)    }^2}\\
	&=\sum_{j=0}^i   \expectation{ \parentheses{ \mathcal{J}_{j,\alpha}(x) \cdot f_j(x) }^2} +\\
	&+ 2\cdot \sum_{j=0}^{i-1}  \sum_{k=j+1}^i   \expectation{ \mathcal{J}_{j,\alpha}(x) \cdot f_j(x) \cdot  \mathcal{J}_{k,\alpha}(x) \cdot f_k(x) } \\
	&= \sum_{j=0}^i    {f_j(x) }^2 \cdot \expectation{ \mathcal{J}_{j,\alpha}(x)} +\\
	&+  2\cdot \sum_{j=0}^{i-1}  \sum_{k=j+1}^i   f_j(x)  \cdot f_k(x) \cdot   \expectation{ \mathcal{J}_{j,\alpha}(x) \cdot   \mathcal{J}_{k,\alpha}(x) } \;.
\end{align*}

For bounding the $\expectation{ \mathcal{J}_{j,\alpha}(x)}$, we can use the fact that $\expectation{ \mathcal{J}_{j,\alpha}(x)}  \leq 2^{i-j}/w^*(T_i)$. For   $\expectation{ \mathcal{J}_{j,\alpha}(x) \cdot   \mathcal{J}_{k,\alpha}(x) }$, we note that the product of the two indicator random variables is $1$ only when both of them are $1$. Having said that, the two random variables are not independent: if $\mathcal{J}_{k,\alpha}(x) =1$, then $ \mathcal{J}_{j,\alpha}(x)=1$ for any $k\geq j$. Namely, if during epoch $k$, $x$ hashed into a counter that was a prefix of $\alpha$, then the same would have happened during epoch $j \leq k$.  We get that:

$$\expectation{ \mathcal{J}_{j,\alpha}(x) \cdot   \mathcal{J}_{k,\alpha}(x) } = \expectation{  \mathcal{J}_{k,\alpha}(x) } \;.$$

\noindent
The whole sums becomes:

\begin{align*}
	\sum_{j=0}^{i-1}  \sum_{k=j+1}^i   f_j(x)  \cdot f_k(x) \cdot   \expectation{ \mathcal{J}_{k,\alpha}(x) } &= \sum_{k=1}^i \sum_{j=0}^{k-1}  f_j(x)  \cdot f_k(x) \cdot   \expectation{ \mathcal{J}_{k,\alpha}(x) }\\
	&= \sum_{k=1}^i   f_k(x) \cdot   \expectation{ \mathcal{J}_{k,\alpha}(x) } \cdot \sum_{j=0}^{k-1}  f_j(x)\\
	&=  \sum_{k=1}^i   f_k(x) \cdot   \expectation{ \mathcal{J}_{k,\alpha}(x) } \cdot f_{\leq k-1}(x)
\end{align*}

where, for simplicity, we denoted $f_{\leq {i}}(x) =  \sum_{j=0}^{i}  f_i(x)$ the sum of frequencies accumulated by $x$ through epochs $0$ to $i$ and set $f_{\leq -1}(x)=0$. Then:

\begin{align*}
	\expectation{\parentheses{\mathcal{C}(x)}^2} &= \sum_{j=0}^i   \parentheses{{f_j(x) }^2  + 2 f_j(x) \cdot f_{\leq j-1}(x)} \cdot \expectation{ \mathcal{J}_{j,\alpha}(x)}  \\
	&=  \sum_{j=0}^i   \parentheses{{f_{\leq j}(x) }^2  - {f_{\leq j-1}(x) }^2} \cdot \expectation{ \mathcal{J}_{j,\alpha}(x)}\;.
\end{align*}


While the term ${f_{\leq j}(x) }^2  - {f_{\leq j-1}(x) }^2$ might seem contrited, we note that $\sum_x{f_{\leq j}(x) }^2  - {f_{\leq j-1}(x) }^2 = F_{j} - F_{j-1}$ denotes the increase in the squared $\ell_2$ norm of the stream vector as we move from epoch $j-1$ to epoch $j$. We thus get that:


\begin{align*}
		\variance{\hat{f}(z)} &= \sum_x \expectation{\parentheses{\mathcal{C}(x)}^2} - (f(z))^2\\
		&= \sum_x \sum_{j=0}^i   \parentheses{{f_{\leq j}(x) }^2  - {f_{\leq j-1}(x) }^2} \cdot \expectation{ \mathcal{J}_{j,\alpha}(x)} -  (f(z))^2\\
		&\leq  \sum_{j=0}^{i}  \parentheses{F_{j} - F_{j-1}} \cdot \frac{2^{i-j}}{w^*(T_i)}\;.
\end{align*}

At this point, it becomes perhaps more natural to understand why we defined the size function to depend on the squared $\ell_2$ norm of the stream vector. We now discuss different bounds we get on the variance depending on the size function.

\medskip\noindent
\textbf{Setting $w^*(F) =  F/\epsilon$ in CS.}   In this case, we start with the original sketch $\mathcal{S}_0$ that allows streams for which their squared $\ell_2$ norm is at most some $F_0$. The size function doubles whenever the squared $\ell_2$ norm doubles. In this case $F_{j+1} = 2 \cdot F_j$ for all $j\geq 0$ and, in general, $F_j = 2^j \cdot F_0$ for all $j\geq 0$. In this case, we get that:

\begin{align*}
	\variance{\hat{f}(z)} &\leq  \sum_{j=0}^{i}  \parentheses{F_{j} - F_{j-1}} \cdot \frac{2^{i-j}}{w^*(T_i)} \\
	&\leq \frac{1}{w^*(T_i)} \cdot \parentheses{  F_0\cdot 2^{i}+ \sum_{j=1}^{i}  \parentheses{2^{j} - 2^{j-1}} \cdot F_0 \cdot 2^{i-j}}\\
	&\leq \frac{1}{w^*(T_i)} \cdot  \parentheses{  F_0\cdot 2^{i}+  \sum_{j=1}^{i}  2^{j-1} \cdot F_0 \cdot 2^{i-j}} \\
		&\leq \frac{1}{w^*(T_i)} \cdot  \parentheses{  F_0\cdot 2^{i}+  \sum_{j=1}^{i}  F_0 \cdot 2^{i-1}}\\
		&\leq  \frac{1}{w^*(T_i)} \cdot  \parentheses{  F_i+  (i-1)\cdot F_i/2} =   \frac{F_i}{w^*(T_i)} \cdot  \frac{i+1}{2} \;.
\end{align*}



\medskip\noindent
\textbf{Setting $w^*(F) =  F^{1-c}/\epsilon$ in CS, for some $0<c<1$.} In this case, the size function doubles more rarely, whenever $F_{i+1} = 2^{1/(1-c)} \cdot F_i$. Since $0<c<1$, note that this implies that $F_{i+1}$ is bigger than $2F_i$. This difference is why we avoid the extra dependency on $i$ in the error term. Formally, we have that $F_i = 2^{i/(1-c)} \cdot F_0$ and thus:

\begin{align*}
	\sum_{j=0}^i  2^{i-j} \cdot  (F_{j+1} - F_j) & =	\sum_{j=0}^i  2^{i-j} \cdot  (2^{(j+1)/(1-c)}\cdot T_0 - 2^{j/(1-c)}\cdot T_0)\\
	&=  \sum_{j=0}^i  2^{i-j} \cdot  2^{j/(1-c)} \cdot T_0 \cdot (2^{1/(1-c)} -1) \\
	&= 2^i \cdot T_0 \cdot (2^{1/(1-c)} -1) \cdot   \sum_{j=0}^i   2^{j/(1-c) - j}   \;.
\end{align*}	









\end{document}
\endinput
%%
%% End of file `sample-sigconf.tex'.
